\chapter{Arithmetic Comparisons and Residual Structure}
\label{v4-arith:comparison}

The six comparison problems of Chapter~\ref{v4-arith:closure} are
\(\mathrm{L}_{\mathrm{ACD}}\) (adelic categorical descent),
F2.c (\(\infty\)-categorical Tate restricted tensor product),
F3.BZN\(^{\mathrm{Ar}}\) (arithmetic Ben-Zvi--Nadler identification),
arithmetic Positselski (co-contra equivalence over \(\Lambda\)),
Koszul--Petersson compatibility, and
VB\(^{*}\) (arithmetic Virasoro positivity).
Each comparison has a specified domain and additional conditions for
extension beyond it. In the positivity problem, three different
objects must be distinguished: a scalar Li coefficient, a transform
of a zero-divisor operator, and a polynomial in a Fock energy operator.
A positive scalar trace does not identify these operators or make each
of their quadratic forms nonnegative.

\section{First Residual Decomposition}
\label{v4-arith:sw:table}

\subsection{F2.c: \texorpdfstring{\(\infty\)}{infinity}-Categorical Tate Restricted Tensor Product}

The formal restricted-tensor construction is proved in
Theorem~\ref{v4-arith:tate:thm:F2c-main}. Its inputs include actual
operadic maps \(\eta_v:\mathrm{Sp}\to\mathcal C_v\) on a cofinite
set of places. These maps form the transition functors between
finite tensors. The colimit is taken over that diagram, without
an infinite tensor of unspecified unit-comparison arrows.

Definition~\ref{v4-arith:sw:def:pointed-En} states the pointing
used here, and Theorem~\ref{v4-arith:sw:conj:F2c-tightened}
records the formal \(E_2\) consequence. The specialization in
Theorem~\ref{v4-arith:w5-h:thm:f2c-e2-transport} has the same
inputs. A non-unit central idempotent does not replace them.
The remaining arithmetic problem is to construct the local
structures, their pointings and any independently required
global geometric comparison.

\subsection{\texorpdfstring{\(\mathrm{L}_{\mathrm{ACD}}\)}{L-ACD}: Adelic Categorical Descent}

\(\mathrm{L}_{\mathrm{ACD}}\) in Chapter~\ref{v4-arith:closure} asserts that
the adelic categorical descent from local automorphic categories to the
global automorphic category holds; but four subconditions are conflated
into the single statement:
(i)~dg-categorical proof that the restricted tensor product with
distinguished units preserves the colimit structure;
(ii)~proof that local Hecke-functoriality together with cross-prime
commutativity implies global Hecke-functoriality on the colimit;
(iii)~sph-finite preservation under the colimit;
(iv)~strong approximation for the simply connected derived group,
together with the determinant idele-class quotient for
\(GL_{2}/\mathbb{Q}\).
Subcondition (iv) is a theorem after this separation: Kneser supplies
the simply connected strong-approximation input for \(SL_{2}\), while
Tate's idelic formalism supplies the determinant quotient. The
conditions (i)--(iii) depend on F2.c and are conjectural.
Thus \(\mathrm{L}_{\mathrm{ACD}}\) splits into the four subconditions
\(\mathrm{L}_{\mathrm{ACD}}.1\)--\(\mathrm{L}_{\mathrm{ACD}}.4\);
see Theorem~\ref{v4-arith:sw:thm:LACD-split}.

\subsection{\texorpdfstring{F3.BZN\(^{\mathrm{Ar}}\)}{F3-BZN-Ar}: Arithmetic Ben-Zvi--Nadler Identification}

F3.BZN\(^{\mathrm{Ar}}\) proposes that the arithmetic analogue of the
Ben-Zvi--Nadler Hochschild-trace identification holds. Three
independent obstructions separate the complex-analytic source theorem
from its arithmetic extension.
(R1)~Ben-Zvi--Nadler's proof over \(\mathbb{C}\) uses holomorphic
factorization on \(S^{1}=B\mathbb{Z}\); the arithmetic analogue requires
locally-constant-with-Adams factorization on \(B\widehat{\mathbb{Z}}\),
and the survival of the Hochschild-trace pairing under this qualitative
change is not established via Gaitsgory--Rozenblyum DAG Vol.~II for
formal-Noetherian targets.
(R2)~The arithmetic free loop space requires a fixed descent topology;
the rigid-analytic topology on the Emerton--Gee stack (arXiv:1908.07185~\S2)
is the canonical choice.
(R3)~The Mazur--Kapranov--Reznikov dictionary (primes as knots,
\(\overline{\mathrm{Spec}\,\mathbb{Z}}\) as 3-manifold) is a heuristic
motivation; the technical content runs through the
Emerton--Gee rigid-analytic construction.
The arithmetic statement uses R2's topology, treats R3 as motivation,
and isolates R1's formal-Noetherian extension as the open
subcondition; see Theorem~\ref{v4-arith:sw:thm:BZN-Ar-split}.

\subsection{Arithmetic Co--Contra Comparison}

An arithmetic co--contra correspondence must specify a coalgebra over
\(\Lambda=\mathbb Z_p[[\Gamma]]\), its continuous dual, and the
categories and resolutions to be compared. The abelian nature of
\(\Gamma\) does not identify an oscillator module with a cyclotomic
Iwasawa module. Nor does the vanishing of a module's \(\mu\)-invariant
imply conilpotence of a coalgebra.
Question~\ref{v4-arith:sw:thm:PosAr-Heisenberg} retains the comparison
with these coefficient requirements.

\subsection{Koszul--Petersson Compatibility: Sub-Claim Structure}

The Koszul--Petersson compatibility asserted in Chapter~\ref{v4-arith:closure}
is the identity
\(\langle\cdot,\cdot\rangle^{\mathrm{Pet}}=\langle\cdot,\cdot\rangle^{\mathrm{Koszul}}\circ(\mathrm{id}\otimes\text{BC-inv.})\)
where BC-inv.\ denotes the Bost--Connes involution on the dual.
The assertion conflates seven independent sub-claims:
(K1)~definition of the arithmetic Koszul pairing via adic transfer of
Vol~I's derived-centre pairing;
(K2)~Hecke-equivariant embedding
\(\mathcal{V}^{\mathrm{prim}}\hookrightarrow L^{2}([GL_{2}])^{\mathrm{sph\text{-}fin}}\);
(K3)~specification of the Bost--Connes anti-involution as the
KMS\({}_{1}\) Tomita--Takesaki modular conjugation after the modular
half-density normalisation;
(K4)~archimedean normalisation;
(K5)~reconciliation of the Bost--Connes KMS\({}_{1}\) pairing with
the local Whittaker/Kirillov pairing at each finite prime \(p\), up
to the declared spherical norm factor
\(m_p(1+p^{-1})L_p(1,\mathrm{Ad})\);
(K6)~agreement of adelic Tamagawa measure with the restricted tensor
product of local Haar measures;
(K7)~confirmation that the sph-finite cut is the natural locus for
the identification.
The explicit model calculations fix the conjugations and local
norms. Their transport to the arithmetic carrier still requires
(K2), the arithmetic identifications in (K3)--(K5), and the adelic
comparison in (K6). A local functional equation does not construct
these maps.
The Heisenberg self-duality and Bost--Connes self-duality are each
properties of their respective objects; the sub-claim structure
exhibits that their pairing requires a separate comparison.
See Theorem~\ref{v4-arith:sw:thm:KP-split}.

\subsection{\texorpdfstring{VB\(^{*}\)}{VB*}: Arithmetic Virasoro Positivity and Spectral Hypotheses}

Scalar positivity \(\lambda_n\ge0\) and operator positivity
\(T_n\succeq0\) are different assertions. Operator positivity means
that \(\langle T_nv,v\rangle\ge0\) for every vector in the specified
quadratic-form domain. It cannot be recovered from the sign of one
trace. This distinction already occurs for bounded self-adjoint
matrices, independently of any issue of moment determinacy.

For example, on \(\mathbb C^2\) with its usual Hermitian inner product,
set \(A=\operatorname{diag}(0,2)\), let \(I\) be the identity, and put
\(S_n=A^n-(A-I)^n\) for \(n\ge1\). Direct diagonal multiplication gives
\[
 S_n=\operatorname{diag}\bigl((-1)^{n+1},2^n-1\bigr),\qquad
 \operatorname{Tr}(S_n)=
 \begin{cases}
  2^n-2,&n\text{ even},\\
  2^n,&n\text{ odd}.
 \end{cases}
\]
Every displayed trace is strictly positive, whereas
\(S_2=\operatorname{diag}(-1,3)\) has quadratic form \(-1\) on the
first coordinate vector. Here \(A\) is bounded, nonnegative and
self-adjoint, and the ordinary trace is positive and faithful.
Its moment measure is \(\delta_0+\delta_2\), where \(\delta_a\)
denotes unit mass at \(a\). This measure is determinate: any positive
measure with the same moments has integral of \(x^2(x-2)^2\) equal
to zero, hence is supported on \(\{0,2\}\); its zeroth and first
moments then determine both masses. Thus even compact spectral
support and moment determinacy do not turn the scalar inequalities
into operator positivity. No identification of these traces with the
arithmetic Li coefficients is being made.

For these particular polynomials there is an exact operator criterion.
Let \(A\) be self-adjoint on a Hilbert space \(\mathcal H\), with
spectral projections \(E_A\), and use the common spectral domain
\[
 \mathcal D_{\mathrm{sp}}
 =\bigcup_{R>0}E_A([-R,R])\mathcal H.
\]
Every power of \(A\) preserves this dense domain. Spectral truncations
converge in the graph norm of every polynomial in \(A\), so it is a
common core. Interpret \(p_n(A)\), where
\(p_n(x)=x^n-(x-1)^n\), first on \(\mathcal D_{\mathrm{sp}}\) and then
by its closed spectral-calculus realization. Then
\[
 p_n(A)\succeq0\quad\text{for all }n\ge1
 \quad\Longleftrightarrow\quad A\succeq\tfrac12 I.
\]
Necessity follows from \(p_2(A)=2A-I\). For sufficiency, odd powers
are strictly increasing on \(\mathbb R\), so \(p_n(x)>0\) for odd
\(n\). If \(n\) is even and \(x\ge\tfrac12\), then
\(|x|\ge|x-1|\), hence \(p_n(x)\ge0\). The spectral calculus gives
the asserted positivity. For an unbounded operator, a statement made
only on a smaller common domain also needs a specified form core
and closure. The scalar traces supply neither that domain nor the
spectral lower bound.

The arithmetic Hilbert--P\'olya hypothesis H-P\(^{\mathrm{Ar}}\)
(Hypothesis~\ref{v4-arith:rh:hyp:HPAr}) supplies a separate spectral object:
\(\Theta_\xi=\tfrac12 I+iH_{\mathrm{HP}}\), with a self-adjoint
ordinate operator, specified domain, determinant divisor and
regularised trace. It does not identify \(H_{\mathrm{HP}}\) with a
Virasoro zero-mode \(L_0^{\mathrm{Ar}}\). Its Li transforms are
\[
 \mathcal L_n=I-(I-\Theta_\xi^{-1})^n,
 \qquad \lambda_n=\operatorname{Tr}_{\mathrm{reg}}(\mathcal L_n),
\]
under the determinant and trace hypotheses of
Proposition~\ref{v4-arith:rh:prop:Li-is-Virasoro}. If the self-adjoint ordinate
operator and the determinant-divisor identity are both assumed,
zero-placement already follows from the spectrum of
\(\Theta_\xi\). The scalar Li route uses the determinant and trace
identities separately.
The forward implication of
Theorem~\ref{v4-arith:rh:thm:VB-implies-F-RH} is therefore positive HP--Li
trace \(\Rightarrow\) Li positivity \(\Rightarrow\) zero-placement.
The operator and trace positivity requirements can be separated
explicitly. Under the self-adjoint H-P\(^{\mathrm{Ar}}\) hypothesis,
put \(U=I-\Theta_\xi^{-1}\). At a real spectral ordinate \(t\),
its scalar value is \((-1/2+it)/(1/2+it)\), of modulus one.
Thus \(U\) is unitary, and
\[
 \operatorname{Re}\mathcal L_n
 :=\frac{\mathcal L_n+\mathcal L_n^*}{2}
 =\frac{(I-U^n)^*(I-U^n)}2\succeq0.
\]
The unsymmetrized \(\mathcal L_n\) need not be self-adjoint.
To pass from this real-part positivity to \(\lambda_n\ge0\), the
regularised trace must be defined, compatible with adjoints and
nonnegative on the relevant real-part cone, with the required
continuity or Li/Weil approximation. These are additional trace
hypotheses in the positive-trace route of
Corollary~\ref{v4-arith:rh:cor:VB-implies-Li}; regularisation alone supplies
no positivity property.

To deduce this implication from VB\(^{*}\), one would also need a
comparison from the primary Fock operators and their positive forms
to these Li transforms and the regularised trace, preserving the
relevant positivity cone. No such comparison is supplied by the
moment argument. Conversely, zero-placement alone does not construct
the primary Virasoro operators, their common invariant domain,
self-adjoint extensions, or this comparison. On the standard Fock
space the zero-energy vacuum prevents positivity of every Virasoro
polynomial. The precise positive spectral sector and the remaining
arithmetic comparison are described below.

\section{Sharp Subconditions}
\label{v4-arith:sw:corrections}

\subsection{Pointed \(E_{n}\)-monoidal Category}

\begin{definition}[pointed \(E_{n}\)-monoidal category]
\label{v4-arith:sw:def:pointed-En}
A \emph{pointed \(E_{n}\)-monoidal stable presentable \(\infty\)-category}
for \(n\geq1\) consists of an \(E_n\)-algebra \(\mathcal C\) in
\(\mathrm{Pr}^{\mathrm L}_{\mathrm{st}}\) and a specified morphism
\[
 \eta:\mathrm{Sp}\longrightarrow\mathcal C
 \quad\text{in}\quad
 \mathrm{Alg}_{E_n}(\mathrm{Pr}^{\mathrm L}_{\mathrm{st}}).
\]
Write \(u=\eta(\mathbb S)\), where \(\mathbb S\) is the sphere.
The notation \((\mathcal C,u,\eta)\) retains this relation.
Compatibility with the nullary operation gives
\(u\simeq\mathbb1_{\mathcal C}\). Thus \(\eta\) is a functor with
all operadic compatibilities, not an arbitrary arrow
\(\mathbb1_{\mathcal C}\to u\) inside \(\mathcal C\).
\end{definition}

\begin{remark}[why pointed]
\label{v4-arith:sw:rem:why-pointed}
The classical restricted direct product specifies local subgroups
and their inclusions. The categorical tensor construction instead
specifies transition functors by tensoring with the maps \(\eta_v\).
Their nullary and higher compatibilities make these functors
operadic morphisms. A chosen object or central idempotent alone
does not do so. The maps into the colimit need not be fully faithful.
\end{remark}

With these pointings, the formal F2.c assertion is the following theorem.

\begin{theorem}[F2.c with specified pointings]
\label{v4-arith:sw:conj:F2c-tightened}
Let \(\mathcal C_v\) be stable presentable \(E_2\)-monoidal categories
whose tensor products preserve colimits separately. On a cofinite
subset \(V^{\mathrm{sph}}\subseteq V\), supply pointings as in
Definition~\ref{v4-arith:sw:def:pointed-En}. Put
\(E=V\setminus V^{\mathrm{sph}}\). Then
\[
 \bigotimes'_v(\mathcal C_v,u_v)
 :=\operatorname*{colim}_{\substack{S\subseteq V\text{ finite}\\E\subseteq S}}
       \bigotimes_{v\in S}\mathcal C_v
\]
is an \(E_2\)-algebra in \(\mathrm{Pr}^{\mathrm L}_{\mathrm{st}}\).

Transitions insert the supplied maps \(\eta_v\).
The finite-stage maps are \(E_2\)-monoidal, with uniqueness
relative to the complete underlying colimit cocone. Singleton
maps require the condition or additional pointings in (M2) of
Theorem~\ref{v4-arith:tate:thm:F2c-main}.
\end{theorem}

\begin{proof}
Tensoring the actual \(E_2\)-algebra maps with identities gives
the finite-stage diagram in
\(\mathrm{Alg}_{E_2}(\mathrm{Pr}^{\mathrm L}_{\mathrm{st}})\).
Its indexing poset is filtered. The categorical tensor product
preserves filtered colimits separately, so HA~3.2.3.1 constructs
its colimit on the underlying categorical colimit.
Theorem~\ref{v4-arith:tate:thm:F2c-main} gives the cocone and
singleton assertions with precisely these hypotheses.
\end{proof}

\subsection{\texorpdfstring{\(\mathrm{L}_{\mathrm{ACD}}\)}{L-ACD} Split into Four Named Subconditions}

\begin{theorem}[\(\mathrm{L}_{\mathrm{ACD}}\) split]
\label{v4-arith:sw:thm:LACD-split}
\(\mathrm{L}_{\mathrm{ACD}}\) is equivalent to
the conjunction of:
\begin{itemize}
\item[(\(\mathrm{L}_{\mathrm{ACD}}.1\))] \emph{Dg-categorical colimit
commutation.} The restricted tensor product of dg-categories
with pointed units commutes with the sph-finite truncation.
\textbf{Form: conjectural}. Formal assembly is supplied by
Theorem~\ref{v4-arith:sw:conj:F2c-tightened}, but it does not prove
compatibility with the specified sph-finite truncation.
\item[(\(\mathrm{L}_{\mathrm{ACD}}.2\))] \emph{Cross-prime Hecke
coherence.} For distinct primes \(p\neq q\), the local Hecke
operators \(T_{p}\) and \(T_{q}\) commute on the restricted tensor
product of local automorphic categories, with associativity and
unit coherences.
\textbf{Form: conjectural}; follows from local Hecke-functoriality
(Fargues--Scholze) plus a symmetric-monoidal coherence check.
\item[(\(\mathrm{L}_{\mathrm{ACD}}.3\))] \emph{Sph-finite preservation.}
The sph-finite subcategory
\(\mathrm{D}^{\mathrm{sph\text{-}fin}}([GL_{2}(\mathbb{Q}_{v})])\) is
preserved under the restricted-tensor colimit: the colimit of
sph-finite categories embeds in the sph-finite subcategory of the
colimit.
\textbf{Form: conjectural}; plausibly automatic given
(\(\mathrm{L}_{\mathrm{ACD}}.1\)).
\item[(\(\mathrm{L}_{\mathrm{ACD}}.4\))] \emph{Strong approximation and determinant quotient.}
The simply connected group \(SL_{2}/\mathbb{Q}\) satisfies strong
approximation away from a nonempty set of places (Kneser 1966;
Platonov--Rapinchuk 1994). The passage from \(SL_{2}\) to \(GL_{2}\)
keeps the determinant idele-class quotient
\(\mathbb{A}_{\mathbb{Q}}^{\times}/\mathbb{Q}^{\times}\), which is
the Tate line rather than a strong-approximation consequence.
\textbf{Form: theorem.}
\end{itemize}
Under this split, \(\mathrm{L}_{\mathrm{ACD}}\) follows from
(\(\mathrm{L}_{\mathrm{ACD}}.1\)) + (\(\mathrm{L}_{\mathrm{ACD}}.2\)) +
(\(\mathrm{L}_{\mathrm{ACD}}.3\)); (\(\mathrm{L}_{\mathrm{ACD}}.4\))
is theorem.
\end{theorem}

\subsection{\texorpdfstring{F3.BZN\(^{\mathrm{Ar}}\)}{F3-BZN-Ar} Split into Three Sub-Obstructions}

\begin{theorem}[Conditional arithmetic trace comparison]
\label{v4-arith:sw:thm:BZN-Ar-split}
Fix the coefficient category, trace construction and completions.
Assume the following comparison data, with compatible circle actions
and the operations required in F3:
\begin{itemize}
\item[(R1)] For the specified spectral category \(\mathcal C_X\),
an actual equivalence from its chosen Hochschild trace complex to
\(\Gamma(LX,\mathcal O_{LX})\). Here the geometric loop object is
the derived self-intersection of the diagonal. An arithmetic or
profinite loop object requires its own comparison.
\item[(R2)] The specified global Selmer-derived target \(X\), with
its coefficient, formal and singular-support conditions, and the
categorical equivalence transporting the arithmetic trace to
that of \(\mathcal C_X\).
\item[(R3)] An actual equivalence from this loop-function complex
to the full Weil/\((\varphi,\Gamma)\) target, preserving the
stated completion and operations. The inertia object
\(\operatorname{Map}(B\widehat{\mathbb Z},X)\) alone does not
supply this equivalence.
\end{itemize}
Then composition gives the F3 trace comparison on the specified
carrier. These are sufficient data; no formal-stack trace theorem
is inferred from a categorical framework alone.
\end{theorem}
\begin{proof}
Transport the trace along the categorical equivalence in (R2),
then apply (R1) and (R3). Their compatibility preserves the specified
circle action, completion and operations.
\end{proof}
\begin{proposition}[The affine loop calculation]
\label{v4-arith:sw:prop:affine-trace-model}
Let \(k\) be a field and \(A\) a smooth finitely generated
commutative \(k\)-algebra. For the small \(k\)-linear category
\(\operatorname{Perf}(A)\) and ordinary derived loops of
\(X=\operatorname{Spec}A\),
\[
 \mathrm{HH}_k(\operatorname{Perf}(A))
 \simeq A\otimes^{\mathbf L}_{A\otimes_k A}A
 \simeq\Gamma(X\mathbin{\times^{\mathbf R}_{X\times_k X}}X,
                   \mathcal O).
\]
\end{proposition}
\begin{proof}
The perfect generator \(A\) identifies the categorical trace with
the trace of the diagonal \(A\)-bimodule. Resolve that bimodule by
the two-sided bar resolution over \(A\otimes_k A\); tensoring
with \(A\) gives the Hochschild chain complex and the first
formula. The derived affine fibre-product rule gives the second.
The cyclic bar construction supplies the circle action. Extending
this calculation to a formal completion or a nilpotent-support
ind-coherent category requires the actual continuity and support
comparison in (R1); smoothness of an arbitrary stack is not such
a comparison.
\end{proof}

\subsection{Arithmetic Positselski with Abelian Domain Restriction}

\begin{question}[an integral co--contra comparison]
\label{v4-arith:sw:thm:PosAr-Heisenberg}
Let \(C\) be a specified complete coalgebra over
\(\Lambda=\mathbb Z_p[[\Gamma]]\).
Construct its co-derived category of continuous comodules and the
contra-derived category of contramodules over its specified continuous dual.
Determine hypotheses on the coefficient modules, resolutions,
conilpotent filtration, and inverse limits that give an equivalence
\[
 \mathrm D^{\mathrm{co}}(C\text{-comod}_{\Lambda})
 \simeq
 \mathrm D^{\mathrm{ctr}}(C^\vee_{\mathrm{cont}}
                         \text{-contramod}_{\Lambda}).
\]
If additional chiral operations are imposed, construct their
compatibility with this ordinary coefficient comparison.
\end{question}

A field-coefficient equivalence cannot be transferred by an
unspecified unital map \(k\to\mathbb Z_p[[\Gamma]]\).
For \(\Gamma\simeq\mathbb Z_p\), the target is
\(\mathbb Z_p[[T]]\), which has characteristic zero and in which
\(p\) is not invertible. A field of positive characteristic cannot
map to it unitally. A field of characteristic zero contains \(1/p\),
whose image would invert \(p\). Thus there is no such field map.
Reduction to \(\mathbb F_p\) and extension to \(\mathbb Q_p\)
are different operations, and an integral comparison must control
both when it uses them.

\begin{question}[the arithmetic Heisenberg coefficient object]
\label{v4-arith:sw:cor:PosAr-Heis-applies}
Construct the coalgebra, coefficient module, topology, and
\(\Gamma\)-action for the proposed arithmetic Heisenberg realization.
If an Iwasawa vanishing theorem is used, identify the resulting
module with one in that theorem's domain and transport its
\(\mu\)-invariant. Prove separately the hypotheses of
Question~\ref{v4-arith:sw:thm:PosAr-Heisenberg}.
\end{question}

Ferrero--Washington concerns the Iwasawa modules attached to abelian
number fields. A trivial \(\Gamma\)-action on an oscillator lattice
does not give the required identification. A co--contra equivalence still leaves the P1 evaluation and defect
identity to construct. The converse in
Theorem~\ref{v4-arith:thm:P1-resolution} also requires defect detection.

\subsection{Koszul--Petersson Compatibility Split into Seven Sub-Claims}

\begin{theorem}[Sufficient pairing comparison data]
\label{v4-arith:sw:thm:KP-split}
The following data give Koszul--Petersson compatibility on the
specified cuspidal core:
\begin{itemize}
\item[(K1)] A defined arithmetic bilinear Koszul pairing on the
chosen completed carrier and a continuous antilinear involution
\(J\), with Hermitian form \(H(v,w)=B(v,Jw)\).
\item[(K2)] An actual Hecke-equivariant arithmetic placement into
the chosen automorphic Hilbert space, with a dense comparison core.
\item[(K3)] Compatibility with local conjugations. On the prime
Gibbs cyclic carrier at \(\beta=1\),
\(J(\mu_{p^k}\Omega)=p^{k/2}\mu_{p^k}^*\Omega\).
This Hilbert--Schmidt carrier is distinct from the weighted
occupation model; their arithmetic placement is supplied data.
\item[(K4)] The archimedean comparison. For the declared
\(K_0\)-Whittaker function, its norm squared is \(\pi/4\), giving
an isometry of the normalized vacuum line. The full Fock comparison
and transport of Segal coefficient conjugation must also be supplied.
\item[(K5)] The finite-place comparison with the declared
multiplicative Haar measure. For normalized unramified tempered
parameters, the spherical Whittaker norm is
\[
 C_p=m_p(1+p^{-1})L_p(1,\mathrm{Ad}),
 \qquad m_p=\operatorname{vol}(\mathbb Z_p^\times).
\]
The normalized line has unit norm; the weighted occupation/radial
and diagonal Hilbert--Schmidt isometries are those of
Chapter~\ref{v4-arith:bc-glue-k5}. Identifying the actual arithmetic
pairing with them is an additional comparison.
\item[(K6)] Compatible normalized reference vectors, completed
restricted tensor products, and an adelic placement preserving
pairings and the specified global measure. Any continued scalar
normalization must be finite and nonzero and applied to both forms.
\item[(K7)] The chosen cuspidal domain and its completed embedding
are specified. No equality on the Eisenstein continuous spectrum is
inferred from these spherical or discrete calculations.
\end{itemize}
Under these data the transported Hermitian Koszul form equals the
Petersson form on the comparison core and hence on its Hilbert
completion wherever both forms are continuous.
\end{theorem}
\begin{proof}
The local model identities are proved in
Chapters~\ref{v4-arith:kp-compat} and~\ref{v4-arith:bc-glue-k5}.
Under the assumed arithmetic intertwiners they are identities of
the arithmetic local forms. Their products agree on normalized
finite tensors by (K6); linearity and continuity extend equality
to the stated completion. The antilinear map in (K1) converts the
bilinear pairing to the required sesquilinear form. This is a
conditional transport theorem, not a construction of (K2) or (K6).
\end{proof}

\subsection{\texorpdfstring{VB\(^{*}\)}{VB*}: Vacuum Obstruction and Positive Spectral Sectors}

The spectral comparison separates the following inputs:
\begin{itemize}
\item H-P\(^{\mathrm{Ar}}\), with the operator, domain, determinant
and trace data of Hypothesis~\ref{v4-arith:rh:hyp:HPAr}; these are additional
to P1, P2 and P3.
\item Positive HP--Li trace, which implies Li positivity and
zero-placement by Corollary~\ref{v4-arith:rh:cor:VB-implies-Li} and
Theorem~\ref{v4-arith:rh:thm:VB-implies-F-RH}, under their stated hypotheses.
\item An arithmetic comparison from a specified positive energy
sector to this trace condition, with the operations preserved on that
sector stated explicitly.
\end{itemize}

The conformal-energy normalization is decisive for the operator
inequalities. Write \(L_0^{\mathrm{Vir}}\) for the standard Fock
energy and \(T_n^{\mathrm{Vir}}=p_n(L_0^{\mathrm{Vir}})\), with
\(p_n(x)=x^n-(x-1)^n\) as above. The vacuum is a primary vector;
the word ``primary'' does not exclude its energy-zero contribution.

\begin{proposition}[vacuum obstruction and maximal positive spectral sector]
\label{v4-arith:sw:conj:VB-converse}
In the normalization of Lemma~\ref{v4-arith:vb:lem:diagonal-action}, the
nonzero vacuum \(\Omega\) satisfies \(L_0^{\mathrm{Vir}}\Omega=0\).
Consequently
\[
 T_n^{\mathrm{Vir}}\Omega=(-1)^{n+1}\Omega,
 \qquad
 \langle T_{2k}^{\mathrm{Vir}}\Omega,\Omega\rangle
 =-\|\Omega\|^2<0\quad(k\ge1).
\]
Thus positivity of every \(T_n^{\mathrm{Vir}}\) fails on the
vacuum-containing Fock space, independently of any hypothesis on
\(\Theta_\xi\).

For a specified self-adjoint energy operator \(A\) on \(\mathcal H\),
with spectral projections \(E_A\), put
\[
 \mathcal H_+=E_A([1/2,\infty))\mathcal H,
 \qquad
 \mathcal D_+=\bigcup_{R\ge1/2}E_A([1/2,R])\mathcal H.
\]
Then \(\mathcal H_+\) is the largest closed reducing subspace on which
every closed spectral-calculus operator \(p_n(A)\) is nonnegative,
and \(\mathcal D_+\) is a common spectral core for their restrictions.
Here reducing means that the orthogonal projection onto the subspace
commutes with every spectral projection of \(A\).
\end{proposition}

\begin{proof}
On the vacuum, substitution of the eigenvalue zero gives
\(p_n(0)=(-1)^{n+1}\), proving the obstruction. On
\([1/2,\infty)\), the polynomial sign calculation above gives
\(p_n\ge0\). Spectral calculus therefore proves positivity on
\(\mathcal H_+\), and bounded spectral truncations give the stated
common core. Conversely, on any closed reducing subspace with this
positivity property, \(p_2(A)=2A-I\succeq0\). Its spectral support
lies in \([1/2,\infty)\), so the subspace is contained in
\(\mathcal H_+\). This proves maximality without a zero-placement
hypothesis.
\end{proof}

For the standard free-boson Fock space, the energies are the
nonnegative integers and the zero-energy space is the vacuum line.
Hence \(\mathcal H_+=\Omega^\perp\). The non-vacuum positivity of
Corollary~\ref{v4-arith:vb:cor:fock-restricted-positivity} therefore extends to
the closed polynomial realizations on this subspace. This is a
statement about the archimedean Fock factor; it does not identify the
full arithmetic energy operator or its domain.

Reducing the energy operator is weaker than preserving a Virasoro
representation. In the \(c=2\) vacuum representation, write
\(L_m\) for the Virasoro modes, so \(L_0=L_0^{\mathrm{Vir}}\) and
\(L_m\Omega=0\) for \(m>0\). The Virasoro relation gives
\[
 L_2L_{-2}\Omega
 =[L_2,L_{-2}]\Omega
 =(4L_0+c/2)\Omega=\Omega.
\]
In particular \(L_{-2}\Omega\ne0\); it has energy two because
\([L_0,L_{-2}]=2L_{-2}\). Thus a non-vacuum vector has a Virasoro
image in the vacuum line. The vacuum complement is not a Virasoro
subrepresentation. A shift \(L_0\mapsto L_0+aI\) with all other
modes fixed is not an alternative normalization of the same
representation: \([L_1,L_{-1}]=2L_0\) forces \(a=0\).

\begin{remark}[arithmetic sector and trace comparison]
\label{v4-arith:sw:rem:VB-converse-finer}
An arithmetic positivity construction must specify a Hilbert sector,
its self-adjoint conformal-energy operator and common core, establish
\(A\succeq I/2\) there, and state which arithmetic and Virasoro
operations preserve that sector. The spectral projection alone
supplies a reducing sector for \(A\), not those additional operations.
A different arithmetic representation must be constructed if the
intended operations require more than the non-vacuum Fock sector.
To imply HP--Li positivity, such a construction must also provide a
comparison with the Li transforms of \(\Theta_\xi\) and the
regularised trace that preserves the required positive forms and
domains. These are separate construction requirements, not
consequences of RH or of moment determinacy.
\end{remark}

\section{First Logical Core}
\label{v4-arith:sw:post-list}

The six residuals from Chapter~\ref{v4-arith:closure} have the
following logical form.

F1 reduces from the monolithic \(\mathrm{L}_{\mathrm{ACD}}\) to the
conjunction \(\mathrm{L}_{\mathrm{ACD}}.1+.2+.3\); subcondition
\(\mathrm{L}_{\mathrm{ACD}}.4\) is the product of strong approximation
for \(SL_{2}/\mathbb{Q}\) and the determinant idele-class quotient for
\(GL_{2}/\mathbb{Q}\); the first is a theorem of Kneser (1966) and
Platonov--Rapinchuk
(1994).

F2.c is Theorem~\ref{v4-arith:sw:conj:F2c-tightened}, with actual
operadic pointings.
Its arithmetic use still requires the local \(E_2\)-structures,
those pointing maps and the comparison with any independently
defined global category.

F3 follows from the compatible R1, R2 and R3 comparison data of
Theorem~\ref{v4-arith:sw:thm:BZN-Ar-split}. The finite-type trace
model proved here is the affine \(\operatorname{Perf}(A)\)
calculation of Proposition~\ref{v4-arith:sw:prop:affine-trace-model}.
The formal trace, global target and full Weil/\((\varphi,\Gamma)\)
equivalences remain actual comparison inputs.

P1 (Theorem~\ref{v4-arith:thm:P1-resolution}) requires the arithmetic
central-defect identity, Mellin evaluation, and defect detection.
Question~\ref{v4-arith:sw:cor:PosAr-Heis-applies} asks separately
for the coefficient object needed to apply an Iwasawa vanishing theorem.

P2 is theorem and unchanged.

P3 follows from the seven specified comparison data in
Theorem~\ref{v4-arith:sw:thm:KP-split}. The local model norms and
conjugations are explicit; the arithmetic placement and completed
adelic pairing comparison remain to be constructed.

For zero-placement, the available implication is positive HP--Li
trace \(\Rightarrow\) Li positivity \(\Rightarrow\) RH
(Theorem~\ref{v4-arith:rh:thm:VB-implies-F-RH}), with the determinant and trace
identities of H-P\(^{\mathrm{Ar}}\). A derivation from VB\(^{*}\)
requires the additional positive comparison described above.

The relevant comparison theories are Positselski (2011) and
Ferrero--Washington (1979) for the arithmetic Positselski abelian
domain; Flath (1979), Weil (1967), and Godement--Jacquet (1972) for
the Koszul--Petersson sub-claims; and Gaitsgory--Rozenblyum DAG
Vol.~II for F3.BZN\(^{\mathrm{Ar}}\). The spectral comparison must
also distinguish scalar traces from operator quadratic forms.

\section{Second Residual Decomposition}
\label{v4-arith:sw:wave2}

The categorical comparisons require higher coherence data in addition
to their scalar identities. The spectral comparison also requires a
separation between a zero-divisor operator and a primary Virasoro
representation, as in Chapter~\ref{v4-arith:polyakov}.

\subsection{F2.c: \texorpdfstring{\(E_{2}\)}{E2}-Algebra Structure on the Unit Map}

The pointing is a morphism \(\eta:\mathrm{Sp}\to\mathcal C\)
in \(\mathrm{Alg}_{E_2}(\mathrm{Pr}^{\mathrm L}_{\mathrm{st}})\).
It is not an object of
\(\mathrm{Alg}_{E_2}(\mathcal C)_{\mathbb1/}\).
Tensoring these functors forms the finite-stage diagram,
including nullary and higher compatibilities.
Uniqueness of its structured colimit is relative to the full
underlying colimit cocone.

Dunn additivity \(E_2\otimes E_2\simeq E_4\) compares two
supplied \(E_2\)-structures with their interchange coherences
to an \(E_4\)-structure. A single \(E_2\)-structure does not
supply that additional data. The inclusion \(E_2\to E_4\)
induces restriction from \(E_4\)-algebras to \(E_2\)-algebras.
Neither this restriction nor the filtered-colimit theorem
constructs a higher structure on an arbitrary \(E_2\)-category.

\subsection{\texorpdfstring{\(\mathrm{L}_{\mathrm{ACD}}\)}{L-ACD}: Archimedean Stratum and Hecke-at-Infinity}

The first-comparison split into \(\mathrm{L}_{\mathrm{ACD}}.1\)--\(.4\)
covers only finite places. At \(v=\infty\), the spherical condition
is \(K_{\infty}\)-finiteness / admissibility in the sense of
Harish-Chandra (1966), not Iwahori-fixedness; the empty-\(S\) base
case collapses to \(D(BK_{\infty})\), not to the trivial category.
The completed Tate character
\(\Lambda_{\zeta}(s)=\Gamma_{\mathbb{R}}(s)\zeta(s)\) forces
any honest \(\mathrm{L}_{\mathrm{ACD}}\) to include a parallel
archimedean factor. The new subcondition is:
\begin{itemize}
\item[(\(\mathrm{L}_{\mathrm{ACD}}.0\))] \emph{Archimedean stratum.}
At \(v=\infty\), the local factor is
\(D((\mathfrak{g},K_{\infty})\text{-mod}^{\mathrm{sph\text{-}fin}})\)
with admissibility (Harish-Chandra) replacing Iwahori-fixedness;
\(\mathrm{L}_{\mathrm{ACD}}.2\) extends to include
\(U(\mathfrak{gl}_{2})^{K_{\infty}}\)-equivariance commuting with
the finite-prime Hecke operators \(T_{p}\); the empty-\(S\) base is
\(\mathrm{Rep}(K_{\infty})\) with
\(\eta:\mathbb{1}\to\mathrm{triv}_{K_{\infty}}\).
\textbf{Form: conjectural.}
\end{itemize}

\subsection{\texorpdfstring{F3.BZN\(^{\mathrm{Ar}}\)}{F3-BZN-Ar}: Formal-Noetherian Extension, Non-Semisimple Correction, and Loop-Space Convention}

The first-comparison subconditions R1--R3 require four refinements.

(D1)~The finite-type model proved in
Proposition~\ref{v4-arith:sw:prop:affine-trace-model} uses
\(\operatorname{Perf}(A)\), a smooth affine \(k\)-scheme and
ordinary derived loops. Denote this precise model by
\(\mathrm{R1.fin\text{-}type}\). The arithmetic
\(\mathrm{R1.formal}\) requirement is an actual comparison for
the chosen completed, support-restricted category on the
Selmer-derived target. It must supply continuity under completion
and the loop-object comparison; neither the phrase ``formal
Noetherian'' nor smoothness of an unspecified stack supplies them.

(D2)~R2 requires an ind-level cotangent-complex calculation on the
global Selmer-derived target. The local smoothness input of
Emerton--Gee \S2 does not by itself supply the global comparison
or its descent compatibility.

(D3)~New subcondition R4: the Iwahori Hecke category
\(\mathrm{Hk}^{\mathrm{Iw}}_{GL_{2},p}\) at wild primes
\(p\in\{2,3\}\) is not semisimple; the Hochschild trace picks up
\(\mathrm{Ext}^{1}\)-contributions invisible to the Ben-Zvi--Nadler
semisimple template (Etingof--Ostrik 2004).
\textbf{Form: conjectural.}

(D4)~R3 requires a derived-internal-Hom convention for profinite
classifying stacks and comparison maps among the cofiltered-limit,
derived internal-Hom and direct ind-constructions after the specified
formal-Noetherian completion. Their agreement is part of the required
comparison.

\subsection{Arithmetic Positselski: Five Supplementary Subconditions}

The coefficient comparison of
Question~\ref{v4-arith:sw:thm:PosAr-Heisenberg} must be supplemented
by further operations in the arithmetic theory.
(C1) requires factorization-coalgebra coherence at each place.
(C2) requires \(d^{2}=0\) for the three-piece arithmetic differential
\(d^{\mathrm{Ar}}=d_{\mathrm{Iw}}+d_{\mathrm{Gersten}}+d_{\mathrm{bar}}^{\mathrm{local}}\)
(the candidate cancellation input is Parshin reciprocity).
(C3) requires Euler-product convergence for the infinite-product structure,
with an actual convergence theorem.
(C4) requires a comparison at the archimedean place.
(C5) requires naturality of the co-contra equivalence in morphisms of
abelian \(\Lambda\)-coalgebras.

C5 is required for any functorial use of a constructed co--contra
comparison. P1 additionally needs the explicit evaluation in
Theorem~\ref{v4-arith:thm:P1-resolution}. Naturality alone supplies
neither its central-defect identity nor its detection condition.

\subsection{Koszul--Petersson Compatibility: Five Higher-Structural Sub-Claims}

The scalar-level decomposition K1--K7 of the first comparison
(Theorem~\ref{v4-arith:sw:thm:KP-split}) conceals five higher-structural
conditions:
(K1a)~the categorical-to-Hilbert descent functor
\(F=\mathrm{RHom}(\mathbb{1},-)\) is not exhibited;
(K3a)~Galois-KMS\({}_{1}\) equivariance requires the Connes--Marcolli
(2004) passage from the Bost--Connes Galois symmetry to the
half-density normalised Tomita--Takesaki conjugation;
(K5a)~defect-pairing commutativity at ramified primes is a potential
anomaly in the Kim arithmetic Chern--Simons defect setup: the open
question is whether the Kim arithmetic-CS defect structure is
anomaly-free at primes where Hecke characters ramify;
(K6a)~the actual adelic measure comparison must be compatible with
any finite nonzero continued \(L(1,\mathrm{Ad}\pi)\) normalization
applied to both pairings;
(K7a)~dimensional-reduction coherence
(\(3d\) adelic \(\to\) \(2d\) archimedean \(\to\) \(1d\) prime-knot)
is stated without proof.
K5a asks for the ramified defect-pairing calculation and its
compatibility with the specified arithmetic comparison maps.

\subsection{\texorpdfstring{VB\(^{*}\)}{VB*}: Zero-Divisor Operators and Primary Virasoro Representations}

H-P\(^{\mathrm{Ar}}\) (Hypothesis~\ref{v4-arith:rh:hyp:HPAr}) concerns the
self-adjoint ordinate operator \(H_{\mathrm{HP}}\) and the normal
operator \(\Theta_\xi=\tfrac12 I+iH_{\mathrm{HP}}\).
A primary Virasoro representation is a separate construction.
Chapter~\ref{v4-arith:polyakov} identifies three obstructions on the
full Hochschild trace class \(\mathcal{V}^{\mathrm{prim}}\):
(i) only the modes \(m\ge-1\) act as global finite sums on the
restricted tensor product;
(ii) the central term \(c_{\mathrm{full}}=\sum_p c^{(p)}\) diverges;
(iii) the scalar \(r_\zeta=2\) is a determinant-rank residue, not a
trace or a global central charge.
Restricting to the primary subspace does not by itself construct
closed Virasoro operators or identify them with the spectral operator.
The relevant distinctions in Chapter~\ref{v4-arith:rh-reformulation} are:
\begin{itemize}
\item Remark~\ref{v4-arith:rh:rem:Virasoro-reconciliation} separates
\(H_{\mathrm{HP}}\) from conformal modes and keeps
\(\mathcal{V}^{\mathrm{prim}}\) as a Hochschild trace class.
\item Conjecture~\ref{v4-arith:rh:conj:primary-Virasoro} asks for closed
operators \(L_m^{\mathrm{PV}}\) on
\(\mathcal{V}^{\mathrm{prim}}_{\mathrm{Heis},\mathrm{prim}}\)
satisfying the Virasoro relations with central charge \(2\).
Its equality with \(r_\zeta=2\) is additional compatibility data.
\item Remark~\ref{v4-arith:rh:rem:spectral-convention} distinguishes the
Fock energy \(L_0^{\mathrm{Vir}}\), Bost--Connes Hamiltonian
\(N_{\mathrm{BC}}\), ordinate operator \(H_{\mathrm{HP}}\), and
zero-divisor operator \(\Theta_\xi\). Self-adjointness of
\(H_{\mathrm{HP}}\) means that its spectral ordinates are real;
\(\Theta_\xi-\tfrac12 I=iH_{\mathrm{HP}}\) is skew-adjoint.
\item Remark~\ref{v4-arith:rh:rem:trivial-zero} uses
\(\xi_{\mathbb R}(s)=\tfrac12s(s-1)\pi^{-s/2}\Gamma(s/2)\zeta(s)\).
The gamma poles cancel the trivial zeros and \(s(s-1)\) removes the
Tate poles. This normalization is distinct from restricting a Fock
representation to a primary subspace.
\end{itemize}
The primary-Virasoro construction and the Weyl-asymptotic comparison
remain separate questions. The latter asks whether the
Riemann--von Mangoldt density
\(N(T)\sim(T/2\pi)(\log(T/2\pi)-1)\) arises in a specified
Berry--Keating \(xp\), Connes adelic, or Meyer--Schwartz operator model.

\section{Resulting Logical Core}
\label{v4-arith:sw:wave2-corrections}

The subcondition structure is as follows.

For F2.c, the pointing in
Definition~\ref{v4-arith:sw:def:pointed-En} is a morphism
\(\mathrm{Sp}\to\mathcal C\) of \(E_2\)-algebras in
\(\mathrm{Pr}^{\mathrm L}_{\mathrm{st}}\). It includes nullary
compatibility. The colimit structure has a contractible space
of choices relative to the entire underlying colimit cocone,
as in Theorem~\ref{v4-arith:sw:conj:F2c-tightened}.

For \(\mathrm{L}_{\mathrm{ACD}}\): new subcondition
\(\mathrm{L}_{\mathrm{ACD}}.0\) (archimedean stratum, as above);
\(\mathrm{L}_{\mathrm{ACD}}.2\) is extended to include
\(U(\mathfrak{gl}_{2})^{K_{\infty}}\)-equivariance.

For F3.BZN\(^{\mathrm{Ar}}\): R1 splits into
\(\mathrm{R1.fin\text{-}type}\) (the stated smooth affine
\(\operatorname{Perf}(A)\) calculation) and \(\mathrm{R1.formal}\)
(the required completed, support-compatible trace comparison); R4 is the non-semisimple
Hecke condition at wild primes \(p\in\{2,3\}\).

For arithmetic co--contra comparison, C1--C5 specify further
operations and naturality. The P1 application also requires the
central-defect identity, Mellin evaluation, and detection condition.

For Koszul--Petersson: five higher-structural sub-claims K1a--K7a
are required; K5a asks whether the ramified defect pairing in the
Kim arithmetic Chern--Simons construction is anomaly-free.

For VB\(^{*}\), the primary-Virasoro construction
(Conjecture~\ref{v4-arith:rh:conj:primary-Virasoro}), its common operator domain,
and its comparison with \(\Theta_\xi\) remain additional data.
The Weyl-asymptotic problem concerns the spectral realization rather
than a Virasoro central-charge computation.

The P1 criterion requires its arithmetic evaluation and defect
detection (Theorem~\ref{v4-arith:thm:P1-resolution}). P2 retains its
stated domain. P3 requires the arithmetic placement, pairings, completions, and
adelic maps of K1--K7 and their higher conditions. The local model
calculations do not supply these comparisons.

The F1--F3 comparisons include
\(\mathrm{L}_{\mathrm{ACD}}.0\), R4, and
\(\mathrm{R1.formal}\). The positive HP--Li trace implication
in Theorem~\ref{v4-arith:rh:thm:VB-implies-F-RH} requires
H-P\(^{\mathrm{Ar}}\) from Hypothesis~\ref{v4-arith:rh:hyp:HPAr}.
The primary Virasoro comparison remains
Conjecture~\ref{v4-arith:rh:conj:primary-Virasoro}.

The sufficient trace route to RH uses the determinant and trace
identities of H-P\(^{\mathrm{Ar}}\) and HP--Li positivity, on the
specified arithmetic Hermitian domain. Arithmetic Positselski
(C1--C5), Koszul--Petersson compatibility (K1--K7, K1a--K7a),
primary-Virasoro reconstruction and Weyl asymptotics are distinct
comparison problems; their labels do not identify their carriers.
A route through operator positivity must first specify its sector and
preserved operations, and then supply the positive trace comparison.
Proposition~\ref{v4-arith:sw:conj:VB-converse} rules out the all-Fock assertion
under the vacuum normalization and constructs the maximal reducing
positivity sector of a given self-adjoint energy operator. Its
comparison with arithmetic operations and the HP--Li trace is the
remaining problem of Remark~\ref{v4-arith:sw:rem:VB-converse-finer}.

The finite trace calculation and the operator dictionary have different
roles. The calculation excludes a general scalar-trace-to-operator
positivity argument. The dictionary states which arithmetic operators
and maps a proof must actually construct. Neither removes the
primary-Virasoro or spectral-realization problems.

The remaining open subconditions are primary-Virasoro reconciliation,
the non-semisimple Hecke condition R4, the defect-pairing anomaly
K5a, and Weyl-asymptotic class.
